\DocumentMetadata{
  pdfversion=2.0,
  pdfstandard=ua-2,
  lang=en-US,
  tagging=on,
  tagging-setup={math/setup=mathml-SE,tabsorder=structure}
}
\documentclass[11pt,letterpaper]{article}
\usepackage[margin=0.68in,includefoot]{geometry}
\usepackage{newcomputermodern}
\usepackage{amsmath,amscd,mathtools}
\usepackage{tikz,adjustbox}
\providecommand{\square}{\mdlgwhtsquare}
\usepackage[hidelinks]{hyperref}
\hypersetup{
  pdftitle={Analysis First-Year Exam, August 2022, Part 2},
  pdfauthor={University of Florida Department of Mathematics},
  pdfsubject={Graduate mathematics examination},
  pdfdisplaydoctitle=true,
  bookmarksopen=true
}
\setcounter{secnumdepth}{0}
\setlength{\parindent}{0pt}
\setlength{\parskip}{0.28em}
\setlength{\emergencystretch}{3em}
\setlength{\leftmargini}{2.15em}
\setlength{\leftmarginii}{2.65em}
\setlength{\leftmarginiii}{3em}
\renewcommand{\labelenumi}{\textbf{\theenumi.}}
\pagestyle{empty}
\raggedbottom
\allowdisplaybreaks
\newenvironment{examproblems}[1]{%
  \begin{enumerate}%
  \setcounter{enumi}{#1}%
  \setlength{\topsep}{0.35em}%
  \setlength{\partopsep}{0pt}%
  \setlength{\itemsep}{0.48em}%
  \setlength{\parsep}{0.18em}%
}{\end{enumerate}}
\newenvironment{examsubparts}{%
  \begin{enumerate}%
  \setlength{\topsep}{0.18em}%
  \setlength{\partopsep}{0pt}%
  \setlength{\itemsep}{0.16em}%
  \setlength{\parsep}{0.1em}%
}{\end{enumerate}}
\newenvironment{examinstructions}{%
  \par\begingroup\itshape
}{%
  \par\endgroup\vspace{0.18em}
}
\makeatletter
\renewcommand\section{\@startsection{section}{1}{0pt}%
  {-1.25ex plus -.25ex minus -.1ex}%
  {0.55ex plus .1ex}%
  {\normalfont\large\bfseries}}
\renewcommand{\@maketitle}{%
  \newpage\null\vskip -1em%
  {\footnotesize\bfseries
    \href{https://www.ufl.edu/}{University of Florida}, \href{https://math.ufl.edu/}{Department of Mathematics}\par}%
  \vskip -0.5em%
  \tagstructbegin{tag=Title}%
    {\LARGE\bfseries\href{https://gma.math.ufl.edu/exams/analysis-first-year/}{Analysis First-Year Exam}, August 2022, Part 2\par}%
  \tagstructend%
  \vskip 0.25em%
  \tagmcbegin{artifact}%
    \hrule height 1pt%
  \tagmcend%
  \vskip 1em%
}
\makeatother
\title{Analysis First-Year Exam, August 2022, Part 2}
\author{}
\date{}
\begin{document}
\maketitle
\thispagestyle{empty}
\begin{examinstructions}
Answer FOUR questions in detail. State carefully any results used without proof.
\end{examinstructions}
\begin{examproblems}{0}
\item
Let \(f\colon [0,1]\to\mathbb{R}\) satisfy the following rules: if~\(t\) is irrational then \(f(t)=0\); if \(t=m/n\) is a rational in lowest terms then \(f(t)=1/n\). Prove that~\(f\) is Riemann integrable and calculate its integral.
\item
Let \(\mathcal{F}\subseteq C(X)\) be equicontinuous and let~\(A\) be the set comprising all points of~\(X\) at which~\(\mathcal{F}\) is bounded; that is, \(a\in A\) exactly when there exists~\(K\) such that \(|f(a)|<K\) whenever \(f\in\mathcal{F}\). Prove that \(A\subseteq X\) is both closed and open.
\item
Let the function \(f:\mathbb{R}\to\mathbb{R}\) be differentiable to all orders and denote its~\(n\)th derivative by \(f_n\). Assume that the sequence \((f_n)_{n=0}^{\infty}\) converges uniformly to the function~\(g\) on~\(\mathbb{R}\) and deduce as much as possible about~\(g\).
\item
Prove that if the real-valued functions~\(f\) and~\(g\) on the same space are measurable, then so are their pointwise sum \(f+g\) and pointwise product \(fg\).
\item
For each \(n\in\mathbb{N}\) let \(f_n:[0,1]\to[0,1]\) be continuous and assume that \(f_n\to0\) pointwise as \(n\to\infty\). Does it follow that 
\[
\int_0^1 f_n(t)\,dt\to0\quad\text{as }n\to\infty?
\]
 Does your answer change if continuous is replaced by Riemann integrable?
\end{examproblems}
\end{document}
